%%
%% winter-lecture24.tex
%% 
%% Made by Alex Nelson <pqnelson@gmail.com>
%% Login   <alex@lisp>
%% 
%% Started on  2026-02-28T08:24:03-0800
%% Last update 2026-02-28T08:24:03-0800
%% 

\lecture[K-Theory]{}

\begin{node}
There is algebraic $K$-theory which is more complicated. We will be
focusing on topological $K$-theory, which is a generalized cohomology
theory in the sense that for each space $X$ there is a $K(X)$ such
that each map $f\colon X\to Y$ between spaces becomes $f^{*}\colon K(Y)\to K(X)$
(represented by the $\Omega$-spectrum $KU$ [resp., $KO$] for complex
[resp., real] $K$-theory). We could take $K(X)=[X,KU]$ to be the
homotopy equivalence classes of maps, but this is not the traditional
(or a useful) way to think of it.

The original definition: we are studying isomorphism classes of vector
bundles over the space $X$ and applying an algebraic construction to
take the ``formal differences'' (to obtain a group). This is analogous
to constructing the integers $\ZZ$ from the natural numbers $\NN$ as
an ordered pair of natural numbers $(m,n)$ interpreted as $m-n$.
\end{node}

\subsection{Vector bundles}

\begin{node}
We can think ofa vector bundle as a special case of fiber bundles, in
the sense that all the fibers are \emph{vector spaces}. This is a
``family'' version of linear algebra, where the family is parametrized
by the base of the fiber bundle. More formally\dots
\end{node}

\begin{definition}
A (Complex) \define{Vector Bundle} $\xi$ over $B$ [assume $B$ is connected]
consists of the following data:
\begin{enumerate}
\item a total space $E$;
\item a continuous map $p\colon E\to B$ called the \define{Projection Map};
\item for each $b\in B$, the structure of a vector space over $\CC$ on
  the fiber $F_{b}=p^{-1}(\{b\})$ (i.e., a scalar multiplication
  $\CC\times F_{b}\to F_{b}$ and vector addition $F_{b}\times F_{b}\to F_{b}$
  and a zero vector)
\end{enumerate}
such that
\begin{itemize}
\item\textsc{Local Triviality Condition}: For each $b\in B$ there
  exists an open neighborhood $U\subset B$ of $b\in U$ and there
  exists a homeomorphism $h\colon U\times\CC^{n}\to p^{-1}(U)$
  suchthat for any $b'\in U$, the map sending
\begin{equation}
\begin{split}
\CC^{n}\to p^{-1}(\{b'\})\\
x\mapsto h(b',x)
\end{split}
\end{equation}
is a vector space isomorphism (``a linear map''); such a pair $(U,b)$
is called a \define{Local Chart} (or a \define{Trivialization Chart}).
\end{itemize}
Such a vector bundle is called a \define{$\CC^{n}$-bundle}.
\end{definition}

\begin{remark}
If we replace $\CC$ by $\RR$, we get the definition of a real vector
bundle, an $\RR^{n}$-bundle. More generally, we could take any field.
\end{remark}

\begin{remark}
There is a slightly different perspective. In the definition, we could
have an abstract linear structure on each fiber such that the
trivialization respects this linear structure.

But another way to think of it is you can think of the linear
structure on each fiber $F_{b}$ is given by local charts. But if we
have two trivialization charts $U_{\alpha}$ and $U_{\beta}$ which
intersect on
\begin{equation}
U_{\alpha\beta}=U_{\alpha}\cap U_{\beta},
\end{equation}
then there is some requirement on the transition map for consistency.

On $U_{\alpha\beta}$ there are two linear structures coming from the
two trivializations
\begin{subequations}
\begin{equation}
h_{\alpha}|_{U_{\alpha\beta}}\colon U_{\alpha\beta}\times\CC^{n}\to p^{-1}(U_{\alpha\beta})
\end{equation}
and
\begin{equation}
h_{\beta}|_{U_{\alpha\beta}}\colon U_{\alpha\beta}\times\CC^{n}\to p^{-1}(U_{\alpha\beta}).
\end{equation}
\end{subequations}
These linear structures should be ``compatible'', in the sense: for
each $b\in U_{\alpha\beta}$, there is a map $\CC^{n}\to\CC^{n}$ via
composing
\begin{equation}
\{b\}\times\CC^{n}\xrightarrow{h_{\beta}}F_{b}\xrightarrow{h_{\alpha}^{-1}}\{b\}\times\CC^{n}
\end{equation}
this should be a linear map, moreover it should be a linear
isomorphism. We define the \define{Transition Map}
\begin{equation}
\begin{split}
\varphi_{\alpha\beta}\colon&U_{\alpha\beta}\to\GL_{n}(\CC)\\
&b\mapsto h^{-1}_{\alpha}|_{\{b\}\times\CC^{n}}\circ h_{\beta}|_{\{b\}\times\CC^{n}}.
\end{split}
\end{equation}
\end{remark}

\begin{example}
We have \define{Trivial Vector Bundles} look like $\CC^{n}\times B$
where $p\colon\CC^{n}\times B\to B$ is just the usual projection map.
\end{example}

\begin{example}
The M\"{o}bius strip 
\begin{equation}
M:=(I\times I)/((0,t)\sim(1,1-t)).
\end{equation}
The projection map $M\to I$ sends $(s,t)\mapsto s$. The fiber is $\RR$
in this case.
\end{example}

\begin{definition}
A \define{Section} of a vector bundle is a continuous map $s\colon B\to E$
such that $p\circ s=\id_{B}$. A section sends $b\in B$ to $s(b)\in F_{b}$.
\end{definition}

\begin{remark}
The \define{Zero Section} is the section which is homeomorphic to the
base $B$.
\end{remark}

\begin{lemma}
A line bundle (i.e., a $\CC^{1}$-bundle) is trivial if and only if it
admits a nowhere zero section (i.e., a section $s\colon B\to E$ such
that for each $b\in B$ we have $s(b)\neq \vec{0}_{b}\in F_{b}$).
\end{lemma}

\begin{proof}
\forwardproof\ Obvious.

\backwardproof\ The nowhere zero section gives a base vector: for each
$x\in E$ with $x=(b,x_{b})$ then $x_{b}=\lambda_{b} s(b)$ for some $\lambda_{b}\in\CC$.
Then we get an isomorphism for
\begin{equation}
\begin{split}
B\times\CC^{1}\to E\\
(b,\lambda)\mapsto(b,\lambda s_{b})
\end{split}
\end{equation}
with inverse
\begin{equation}
\begin{split}
E\to B\times\CC^{1}\\
(b,x_{b})\mapsto(b,\lambda_{b}).
\end{split}
\end{equation}
Hence the claim.
\end{proof}

\begin{lemma}
A $\CC^{n}$-bundle is trivial if and only if there exists $n$ sections
$s_{1}$, \dots, $s_{n}$ such that for any $b\in B$ the vectors
$s_{1}(b)$, \dots, $s_{n}(b)$ are linearly independent in the fiber $F_{b}$.
\end{lemma}

\begin{proof}
Same as the $\CC^{1}$ case.
\end{proof}

\begin{example}[Tangent bundles]
Given a smooth manifold $M$, then its tangent bundle
\begin{equation}
TM=\{(b,v)\in M\times T_{b}M\}
\end{equation}
with projection map
\begin{equation}
\begin{split}
p\colon & TM\to M\\
& (b,v)\mapsto b
\end{split}
\end{equation}
The local charts of $TM$ are given by local charts of $M$.
\end{example}

\begin{definition}\label{defn:pullback-bundle}
Given a vector bundle $p\colon E\to B$ and a continuous map $f\colon A\to B$,
the \define{Pullback Bundle} $p^{*}\colon f^{*}E\to A$ consists of
\begin{equation}
f^{*}E=\{(a,x)\in A\times E\mid f(a)=p(x)\}.
\end{equation}
Or, diagrammatically,
\begin{equation}
\vcenter{\xymatrix{f^{*}E\ar[d]_{p^{*}}\ar[r]\ar@{}[dr]|>{\txt{$\ulcorner$}} & E\ar[d]^{p}\\
A\ar[r]^{f} & B}}
\end{equation}
The local chart of pullback bundle is obtained by pulling back local
charts $U_{\alpha}$ of $p\colon E\to B$ to $f^{-1}(U_{\alpha})$.
\end{definition}

\begin{remark}
Pullbacks are more natural than pushforwards because the dimension of
the fiber for pullback is constant, but for pushforward this is not
necessarily the case.
\end{remark}

\begin{slogan}
Anything you can do in linear algebra has a vector bundle version.
\end{slogan}

\begin{recall}
In linear algebra, we had the Hermitian bilinear form
\begin{equation}
\langle-,-\rangle\colon\CC^{n}\times\CC^{n}\to\CC
\end{equation}
such that
\begin{enumerate}
\item $\langle u,v_{1}+v_{2}\rangle=\langle u,v_{1}\rangle+\langle u,v_{2}\rangle$
for all $u,v_{1},v_{2}\in\CC^{n}$
\item it is conjugate-linear in one of the slots $\langle\alpha u,\beta v\rangle=\overline{\alpha}\beta\langle u,v\rangle$
\item $\langle u,v\rangle=\overline{\langle v,u\rangle}$.
\end{enumerate}
We can also define the induced \define{Norm} $\|v\|=+\sqrt{\langle v,v\rangle}$.
\end{recall}

\begin{definition}
A \define{Hermitian Bilinear Form} (resp., a norm) on a $\CC^{n}$-bundle is a
Hermitian bilinear form (resp., a norm) on each fiber which varies continuously as we
move a long $B$.

A norm is a continuous map $\|-\|\colon E\to\RR$ such that for each
fiber $F_{b}$ is restriction on $F_{b}$ is a norm.
\end{definition}

\begin{definition}
A $\CC^{n}$-bundle with a Hermitian bilinear form is called a
\define{Hermitian Bundle}.
\end{definition}

\begin{lemma}
If the base $B$ of a $\CC^{n}$-bundle is paracompact (or more
generally, compact and Hausdorff), then there exists a Hermitian
bilinear form on the bundle.
\end{lemma}

Without loss of generality, we can assume such things exist for the
rest of this course.

\begin{node}
In linear algebra, we had a notion of a ``subspace''. For vector
bundles, the analogous concept is a ``subbundle''.
\end{node}

\begin{definition}
Let $\xi$ and $\eta$ be two vector bundles such that the total space
of $E(\xi)$ contains $E(\eta)$ as a subspace $E(\xi)\supset E(\eta)$.
Then we say $\eta$ is a \define{Subbundle} of $\xi$ if for each $b\in B$
we have $F_{b}(\eta)$ is a subspace of $F_{b}(\xi)$.
\end{definition}

\begin{node}
In linear algebra, we have a notion of ``direct sums'' of vector
spaces. For vector bundles, the counterpart is the ``Whitney sum''.
\end{node}

\begin{definition}
Given two vector bundles $\xi_{1}$ with $p_{1}\colon E_{1}\to B_{1}$
and $\xi_{2}$ with $P_{2}\colon E_{2}\to B_{2}$, the \define{Cartesian Product}
$\xi_{1}\times\xi_{2}$ is the vector bundle $p_{1}\times p_{2}\colon E_{1}\times E_{2}\to B_{1}\times B_{2}$.
\end{definition}

\begin{definition}
Given two vector bundles over the same base space $\xi_{1}$ with
$p_{1}\colon E_{1}\to B$ and $\xi_{2}$ with $p_{2}\colon E_{2}\to B$,
we define the \define{Whitney Sum} $\xi_{1}\oplus\xi_{2}$ to be the
pullback bundle (\S\ref{defn:pullback-bundle}) $\Delta^{*}(\xi_{1}\times\xi_{2})$ where
$\Delta\colon B\to B\times B$ is the diagonal function sending $b\mapsto(b,b)$.
\end{definition}